\documentclass[11pt]{article}
\usepackage[a4paper,margin=21mm]{geometry}
\usepackage{fontspec}
\setmainfont{Latin Modern Roman}
\newfontfamily\CJKfont{Noto Serif CJK SC}
\XeTeXlinebreaklocale "zh"
\XeTeXlinebreakskip=0pt plus 1pt
\usepackage{amsmath,amssymb,amsthm,amscd,graphicx,color}
\usepackage{xurl}
\usepackage[unicode,hidelinks]{hyperref}
\allowdisplaybreaks
\setlength\emergencystretch{3em}
\numberwithin{equation}{section}
\newtheorem{Theorem}{Theorem}[section]
\newtheorem*{Theorem*}{Theorem}
\newtheorem{Corollary}[Theorem]{Corollary}
\newtheorem{Lemma}[Theorem]{Lemma}
\newtheorem{Proposition}[Theorem]{Proposition}
 { \theoremstyle{definition}
\newtheorem{Definition}[Theorem]{Definition}
\newtheorem{Note}[Theorem]{Note}
\newtheorem{Example}[Theorem]{Example}
\newtheorem{Remark}[Theorem]{Remark} }

\def\N{\mathbb{N}}
\def\C{\mathbb{C}}% complex nember field
\def\R{\mathbb{R}}% real number field
\def\Q{\mathbb{Q}}% rational number field
\def\Z{\mathbb{Z}}% ring of integers
\def\A{\mathbb{A}}% the ring of adeles
\def\F{\mathbb{F}}% finite field
\def\H{\mathbb{H}}% Hamilton's quaternion field

%---------------Algebraic groups------------------------
\def\GL{\mathrm{GL}}
\def\PGL{\mathrm{PGL}}
\def\SL{\mathrm{SL}}
\def\PSL{\mathrm{PSL}}
%\def\O{\textnormal{O}}
\def\SO{\mathrm{SO}}
\def\GSO{\mathrm{GSO}}
\def\U{\textnormal{U}}
\def\SU{\textnormal{SU}}
\def\Sp{\mathrm{Sp}}
\def\GSp{\mathrm{GSp}}
\def\M{\mathrm{M}}% matrix space

\def\gl{\mathfrak{gl}}
\def\sl{\mathfrak{sl}}
\def\so{\mathfrak{so}}

%--------------symbols----------------------
\def\t{\,{}^t}% transpose

\def\d{\,\mathrm{d}}% measure

\def\idots{\reflectbox{$\ddots$}}% anti-diagonal dots

\def\Cc{C_c^{\infty}}% smooth functions

\def\ds{\displaystyle}% displaystyle

\def\bs{\backslash}% backslash

\def\t{\,{}^t}% transpose

\def\idots{\reflectbox{$\ddots$}}% anti-diagonal dots

\def\tbe{\textcolor{red}{\bf To Be Edited.}}

\DeclareMathOperator{\diag}{diag} % diagonal matrix

\DeclareMathOperator{\supp}{supp}% support

\DeclareMathOperator{\ord}{ord}% order

\DeclareMathOperator{\vol}{vol}% volume

\DeclareMathOperator{\val}{val}% valuation

\DeclareMathOperator{\mat}{Mat}% matrix space

\DeclareMathOperator{\rank}{rank}% rank

\DeclareMathOperator{\Gal}{Gal}% Galois group

\DeclareMathOperator{\Res}{Res}% residue / scalar restriction

\DeclareMathOperator{\Ind}{Ind}% induction

\DeclareMathOperator{\ind}{ind}% compact induction

\DeclareMathOperator{\Hom}{Hom}% space of homomorphisms

\DeclareMathOperator{\End}{End}% space of endomorphisms

\DeclareMathOperator{\Aut}{Aut}% group of automorphisms

\DeclareMathOperator{\Lie}{Lie}% Lie algebra

\DeclareMathOperator{\Inn}{Inn}% group of inner automorphisms

\DeclareMathOperator{\Out}{Out}% group of outer automorphisms

\DeclareMathOperator{\Cent}{Cent}% centralizer

\DeclareMathOperator{\Irr}{Irr}% set of irreducible representations

\DeclareMathOperator{\sgn}{sgn}% the sign character

\DeclareMathOperator{\re}{Re}% real part

\DeclareMathOperator{\im}{Im}% imaginary part


\DeclareMathOperator{\Ad}{Ad}% adjoint action
\DeclareMathOperator{\ad}{ad}% adjoint action
\DeclareMathOperator{\Kfin}{\text{$K$-fin}}% the space of K-finite vectors
\DeclareMathOperator{\Ker}{Ker}% kernel
\DeclareMathOperator{\Image}{Im}% image
\DeclareMathOperator{\triv}{triv}% the trivial representation
\DeclareMathOperator{\tr}{tr}% trace
\DeclareMathOperator{\id}{id}% identity operator
\DeclareMathOperator{\Sym}{Sym}% symmetrization
\DeclareMathOperator{\cpt}{cpt}% compact
\DeclareMathOperator{\semi}{ss}% semisimple

%------------fonts------------------------
%<<Roman>>
\def\a{\alpha}
\def\b{\beta}
\def\c{\gamma}
\def\d{\delta}
\def\e{\varepsilon}
\def\f{\varphi}
\def\g{\psi}
\def\h{\hbar}
%\def\i{\mbox{\raisebox{.5ex}{$\chi$}}}
\def\k{\kappa}
\def\l{\lambda}
\def\m{\mu}
\def\n{\nu}
%\def\o{\omega}
\def\ro{\rho}
\def\s{\sigma}
\def\t{\tau}
\def\x{\xi}
\def\y{\eta}
\def\z{\zeta}
\def\th{\theta}
\def\pa{\partial}
\newcommand{\jap}[1]{\langle #1 \rangle}
%<<mathfrak>>
\newcommand{\fa}{\mathfrak{a}}
\newcommand{\fb}{\mathfrak{b}}
\newcommand{\fc}{\mathfrak{c}}
\newcommand{\fe}{\mathfrak{e}}
\newcommand{\ff}{\mathfrak{f}}
\newcommand{\fg}{\mathfrak{g}}
\newcommand{\fh}{\mathfrak{h}}
\newcommand{\fii}{\mathfrak{i}}
\newcommand{\fj}{\mathfrak{j}}
\newcommand{\fk}{\mathfrak{k}}
\newcommand{\fl}{\mathfrak{l}}
\newcommand{\fm}{\mathfrak{m}}
\newcommand{\fn}{\mathfrak{n}}
\newcommand{\fo}{\mathfrak{o}}
\newcommand{\fp}{\mathfrak{p}}
\newcommand{\fq}{\mathfrak{q}}
\newcommand{\fr}{\mathfrak{r}}
\newcommand{\fs}{\mathfrak{s}}
\newcommand{\fS}{\mathfrak{S}}
\newcommand{\ft}{\mathfrak{t}}
\newcommand{\fu}{\mathfrak{u}}
\newcommand{\fv}{\mathfrak{v}}
\newcommand{\fx}{\mathfrak{x}}
\newcommand{\fy}{\mathfrak{y}}
\newcommand{\fz}{\mathfrak{z}}
\newcommand{\fI}{\mathfrak{I}}
\newcommand{\fJ}{\mathfrak{J}}
\newcommand{\fK}{\mathfrak{K}}
\newcommand{\fL}{\mathfrak{L}}
\newcommand{\fP}{\mathfrak{P}}
\newcommand{\fT}{\mathfrak{T}}
\newcommand{\fW}{\mathfrak{W}}
\newcommand{\fX}{\mathfrak{X}}

%<< mathbb>>
%\newcommand{\bA}{\mathbb{A}}
\newcommand{\bB}{\mathbb{B}}
%\newcommand{\bC}{\mathbb{C}}
\newcommand{\bD}{\mathbb{D}}
\newcommand{\bE}{\mathbb{E}}
\newcommand{\bF}{\mathbb{F}}
\newcommand{\bG}{\mathbb{G}}
%\newcommand{\bH}{\mathbb{H}}
\newcommand{\bI}{\mathbb{I}}
\newcommand{\bJ}{\mathbb{J}}
\newcommand{\bK}{\mathbb{K}}
\newcommand{\bL}{\mathbb{L}}
\newcommand{\bM}{\mathbb{M}}
%\newcommand{\bN}{\mathbb{N}}
\newcommand{\bO}{\mathbb{O}}
\newcommand{\bP}{\mathbb{P}}
%\newcommand{\bQ}{\mathbb{Q}}
%\newcommand{\bR}{\mathbb{R}}
\newcommand{\bS}{\mathbb{S}}
\newcommand{\bT}{\mathbb{T}}
\newcommand{\bU}{\mathbb{U}}
\newcommand{\bV}{\mathbb{V}}
\newcommand{\bW}{\mathbb{W}}
\newcommand{\bX}{\mathbb{X}}
\newcommand{\bY}{\mathbb{Y}}
%\newcommand{\bZ}{\mathbb{Z}}

%<<mathcal>>
\newcommand{\cA}{\mathcal{A}}
\newcommand{\cB}{\mathcal{B}}
\newcommand{\cC}{\mathcal{C}}
\newcommand{\cD}{\mathcal{D}}
\newcommand{\cE}{\mathcal{E}}
\newcommand{\cF}{\mathcal{F}}
\newcommand{\cG}{\mathcal{G}}
\newcommand{\cH}{\mathcal{H}}
\newcommand{\cI}{\mathcal{I}}
\newcommand{\cJ}{\mathcal{J}}
\newcommand{\cK}{\mathcal{K}}
\newcommand{\cL}{\mathcal{L}}
\newcommand{\cM}{\mathcal{M}}
\newcommand{\cN}{\mathcal{N}}
\newcommand{\cO}{\mathcal{O}}
\newcommand{\cP}{\mathcal{P}}
\newcommand{\cQ}{\mathcal{Q}}
\newcommand{\cR}{\mathcal{R}}
\newcommand{\cS}{\mathcal{S}}
\newcommand{\cT}{\mathcal{T}}
\newcommand{\cU}{\mathcal{U}}
\newcommand{\cV}{\mathcal{V}}
\newcommand{\cW}{\mathcal{W}}
\newcommand{\cX}{\mathcal{X}}
\newcommand{\cY}{\mathcal{Y}}
\newcommand{\cZ}{\mathcal{Z}}


\newcommand{\sI}{\mathscr{I}}

\newcommand{\bfe}{\mathbf{e}}
\newcommand{\bfh}{\mathbf{h}}
\newcommand{\bfi}{\mathbf{i}}
\newcommand{\bfj}{\mathbf{j}}
\newcommand{\bfk}{\mathbf{k}}
\newcommand{\bfI}{\mathbf{I}}
\newcommand{\bfJ}{\mathbf{J}}
\newcommand{\bfK}{\mathbf{K}}


\begin{document}
\begin{center}\Large Uniform stationary phase with a moving critical point\end{center}
\noindent\textbf{Stable ID:} 553\quad\textbf{Author:} Kouichi Taira; Hiroyoshi Tamori\par
\noindent\textbf{Work:} Strichartz Estimates for the (k,a)-Generalized Laguerre Operators\par
\noindent\textbf{Source:} \url{https://sigma-journal.com/2025/014/}\par
\noindent\textbf{Version:} Official SIGMA 21 (2025) PDF and native TeX; acquired 2026-09-30; exact bytes pinned by SHA256\par
\noindent\textbf{Rights:} CC BY 4.0. Authors retain copyright. \url{https://creativecommons.org/licenses/by/4.0/}. Reuse and typesetting adaptation; original language and mathematical content preserved.\par
\noindent\textbf{Difficulty (prior selection preserved):} {\CJKfont H2: The saddle approaches the vanishing-amplitude endpoint at a parameter-dependent scale. The proof must choose a cutoff, rescale by the moving saddle, and control three regions uniformly in both parameters; a fixed-parameter stationary-phase application does not give the stated extra decay.}\par
\medskip\noindent\textbf{Editorial boundary note (not author text):} Complete original Proposition 2.4, its integral and parameter conventions, and full cutoff/two-regime/rescaled-three-region proof. The actual Leibniz bound and singular-amplitude nonstationary proof are retained, with the standard parameter-uniform stationary-phase statement and author bibliographic references. This is a general phase/amplitude theorem, not merely substitution into a fixed-parameter stationary-phase result. No downstream dispersive estimate is added.\par
\medskip\hrule\medskip\noindent\textbf{Author text begins below.}\par
\setcounter{section}{1}\setcounter{subsection}{4}
% Original source lines 506--513
\subsection{Notation}
Let
\begin{align*}
\mathbb{N}=\{0,1,2,\hdots\},\qquad \mathbb{N}^*=\{1,2,3,\hdots\}, \qquad \R_{+}=(0,\infty),\qquad \R_-=(-\infty,0).
\end{align*}
For a parameter $i\in I$ and $A_i,B_i$, we write $A_i\lesssim B_i$ for $i\in I$ if there is a constant $C>0$ independent of $i\in I$ such that $A_i\leq CB_i$. We denote $A_i\sim B_i$ if both $A_i\lesssim B_i$ and $B_i\lesssim A_i$ hold for $i\in I$.

\section{Preliminary}

\setcounter{subsection}{3}\setcounter{Theorem}{3}\subsection*{Author moving-critical-point estimate}
% Original source lines 576--695
We consider the following integral with parameters~$\l$,~$\f$:
\begin{align*}
I(\l,\f):=\int_{\R}{\rm e}^{{\rm i}\l S(\m,\f)}\c(\m,\l,\f){\rm d}\m
\end{align*}
and its decay rate with respect to $\l\gg 1$

If we freeze the parameter $\f$, then the stationary phase theorem just implies that if $\m\mapsto S(\m,\f)$ is a Morse function (that is, all critical points are non-degenerate in the sense that $\pa_{\m}^2S\neq 0$ there), the optimal decay rate of $I$ is \smash{$\l^{-\frac{1}{2}}$}. To obtain an improved decay, here we consider the following two scenarios:
\begin{itemize}\itemsep=0pt
\item If we assume that $\c$ vanishes with order $2\n$ at all the critical points of $S$ in addition, then the optimal decay rate is improved to be \smash{$\l^{-\frac{1}{2}-\n}$}.
\item If the symbol $\c$ itself decays like $(1+\l)^{-\n}$, then the decay order becomes \smash{$\l^{-\frac{1}{2}-\n}$}.
\end{itemize}
The following proposition interpolates the above two situations in a uniform way with respect to the parameter $\f$.
A typical example of such a phase function is $S(\m,\f)=(\m-\f)^2$. This simple case could simplify the reading of the proof below.

\begin{Proposition}\label{Stphasemovecrit}
Let $\n\geq 0$ and $c, C>0$. Suppose that $S\in C^{\infty}\bigl([0,1]^2;\R\bigr)$ and $\c(\cdot,\l,\cdot)\in C^{\infty}(\R\times [0,1])$ satisfy
\begin{align*}
&\supp \c(\cdot,\l,\f)\subset [0,1],\qquad |\pa_{\m}^{\a}\c(\m,\l,\f)|\leq C_{\a}\m^{2\n-\a}(1+\l \f\m)^{-\n},\\
&
\big|\pa_{\m}^2S(\m,\f)\big|\geq c,\qquad |\pa_{\m}\pa_{\f}S(\m,\f)|\geq C%\leq C
\end{align*}
uniformly in $\m\in [0,1]$, $\f\in [0,1]$ and $\l\gg 1$. Suppose that $S(\cdot,\f)$ has a unique critical point~${\m_0(\f)}$ which is smooth with respect to $\f$ and that $\m_0(0)=0$. Then
 \smash{$|I(\l,\f)|\lesssim \l^{-\n-\frac{1}{2}}$} uniformly in $\l\gg 1$ and $\f\in [0,1]$.
\end{Proposition}

\begin{Remark}\quad
\begin{itemize}\itemsep=0pt
\item[(1)] Without the assumption $\m_0(0)=0$, we obtain \smash{$I(\l,\f)\lesssim \l^{-\frac{1}{2}}$} only. The additional decay~$\l^{-\n}$ comes from the factor $(1+\l \f\m)^{-\n}$ in the assumption of $\c$ and the fact that $\m_0(0)=0$ as is mentioned before.

\item[(2)] We observe that the critical point $\m_0(\f)$ of $S$ is close to $0$ if $\f\sim 0$. There we take advantage of the assumption that the symbol $\c$ vanishes at zero with order $2\n$ like $|\c(\m,\l,\f)|\lesssim \m^{2\n}$. On the other hand, if $\f$ is away from zero, then the critical point $\m_0(\f)$ is also away from zero. In this case, $\c$ itself decays like $\l^{-\n}$ near the critical point (due to $|\f|\gtrsim 1$ and~${|\m_0(\f)|\gtrsim 1}$), which makes us to prove the improved decay. The difficulty here is to deal with the case where $\f$ is small enough but depending on $\l$.

\item[(3)]
The critical point of $S(\cdot,\f)$ is unique due to the condition $\big|\pa_{\m}^2S(\m,\f)\big|\geq c$ in this case.
\end{itemize}
\end{Remark}

\begin{proof}[Proof of Proposition~\ref{Stphasemovecrit}]
We may assume \smash{$\supp \c(\cdot,\l,\f)\subset \bigl[\l^{-\frac{1}{2}},1\bigr]$}. In fact, we set
\begin{align*}
\c'(\m,\l,\f)=\chi\bigl(\l^{\frac{1}{2}}\m\bigr)\c(\m,\l,\f),
\end{align*}
where $\chi\in C^{\infty}([0,\infty),\R)$ satisfies $\chi(\m)=1$ on $\m\leq 1$ and $\chi(\m)=0$ for $\m\geq 2$. Using the bound~${|\c(\m,\l,\f)|\lesssim \m^{2\n}}$, we have
\begin{align*}
\left|\int_{\R}{\rm e}^{{\rm i}\l S(\m,\f)}\c'(\m,\l,\f){\rm d}\m\right|\lesssim \l^{-\n}\left|\int_{\R}\chi\bigl(\l^{\frac{1}{2}}\m\bigr){\rm d}\m\right|\lesssim \l^{-\n-\frac{1}{2}}.
\end{align*}
Hence, we can replace $\c$ by $\c-\c'$, where we note that $\c-\c'$ satisfies
\begin{align*}
\supp (\c-\c')(\cdot,\l,\f)\subset\bigl[\l^{-\frac{1}{2}},1\bigr],\qquad |\pa_{\m}^{\a}(\c-\c')(\m,\l,\f)|\leq C_{\a}\m^{2\n-\a}(1+\l \f\m)^{-\n}.
\end{align*}
In the following, we assume \smash{$\supp \c(\cdot,\l,\f)\subset \bigl[\l^{-\frac{1}{2}},1\bigr]$}.

 $(i)$ First, we consider the case $\f\in [0,1]$ satisfying \smash{$\m_0(\f)\leq 2^{-1}\l^{-\frac{1}{2}}$}.
Since the signature of~$\pa_{\m}^2S(\m,\f)$ does not change for \smash{$\m\in \bigl[\l^{-\frac{1}{2}},1\bigr]$} and $\f\in [0,1]$, we have
\[
|\pa_{\m}S(\m,\f)|=\left|\int_{\m_0(\f)}^{\m}\pa_{\m}^2S(\m',\f){\rm d}\m'\right|=\int_{\m_0(\f)}^{\m}\big|\pa_{\m}^2S(\m',\f)\big|{\rm d}\m'\geq c(\m-\m_0(\f))\geq 2^{-1}c\m
\]
 for \smash{$\m\in \bigl[\l^{-\frac{1}{2}},1\bigr]$}. By integrating by parts and using \smash{$\supp \c(\cdot,\l,\f)\subset \bigl[\l^{-\frac{1}{2}},1\bigr]$}, we have
\begin{align*}
|I(\l,\f)|=\left|(-{\rm i}\l)^{-N}\int_{\R}{\rm e}^{{\rm i}\l S(\m,\f)} L^N\c(\m,\l,\f){\rm d}\m\right|\leq \l^{-N}\int_{\l^{-\frac{1}{2}}}^1|L^N\c(\m,\l,\f)|{\rm d}\m,
\end{align*}
where $L=\pa_{\m}\circ (\pa_{\m}S)(\m,\f)^{-1}$ and we take $N>\n+1$. We observe from the assumption and Lemma \ref{Leibnizrule} that $\big|L^N\c(\m,\l,\f)\big|\lesssim \m^{2\n-2N}$, which leads to
\begin{align*}
|I(\l,\f)|\lesssim \l^{-N}\int_{\l^{-\frac{1}{2}}}^1\m^{2\n-2N}{\rm d}\m\lesssim \l^{-\n-\frac{1}{2}}.
\end{align*}

 $(ii)$ Next, we consider the case $\f\in [0,1]$ satisfying \smash{$\m_0(\f)\in \bigl[ 2^{-1}\l^{-\frac{1}{2}}, 1\bigr]$}.
Differentiating $(\pa_{\m}S)(\m_0(\f),\f)=0$ with respect to $\f$, we have
\begin{align*}
|\m_0'(\f)|=\bigl|-(\pa_{\m}\pa_{\f}S)(\m_0(\f),\f)/\pa_{\m}^2S(\m_0(\f),\f)\bigr|\sim 1,
\end{align*}
and hence $\m_0(\f)\sim \f$ for $\f\in [0,1]$ due to $\m_0(0)=0$. Now the assumption \smash{$\m_0(\f)\geq 2^{-1}\l^{-\frac{1}{2}}$} yields \smash{$\f\gtrsim \l^{-\frac{1}{2}}$}.

By scaling, we have
\begin{align*}
I(\l,\f)=\int_{\R}{\rm e}^{{\rm i}\l S(\m,\f)}\c(\m,\l,\f){\rm d}\m=\m_0(\f)\l^{-\n}\int_{\R}{\rm e}^{{\rm i}\l \m_0(\f)^2S_{\f}(\m)}\c_{\f}(\m,\l){\rm d}\m,
\end{align*}
where we set
\begin{align*}
S_{\f}(\m)=\m_0(\f)^{-2}S(\m_0(\f)\m,\f),\qquad \c_{\f}(\m,\l)=\l^{\n}\c(\m_0(\f)\m,\l,\f).
\end{align*}
We note that $\m_0(\f)$ is a unique critical point of $S(\m,\f)$.
Let $\g_1,\g_2,\g_3\in C^{\infty}(\R;[0,1])$ such that~${\g_1+\g_2+\g_3=1}$, $\supp \g_1\subset \bigl(-\infty,\frac{3}{4}\bigr]$, $\supp \g_2\subset \bigl[\frac{1}{2},\frac{3}{2}\bigr]$ and $\supp \g_3\subset \bigl[\frac{5}{4},\infty\bigr)$. We write
\begin{align*}
I(\l,\f)
=\m_0(\f)\l^{-\n}\sum_{j=1}^3\int_{\R}{\rm e}^{{\rm i}\l \m_0(\f)^2S_{\f}(\m)}\c_{\f}(\m,\l)\g_j(\m){\rm d}\m
=\m_0(\f)\l^{-\n}\sum_{j=1}^3I_j(\l,\f).
\end{align*}

First, we deal with $I_2(\l,\f)$.
Now we see $|\pa_{\m}S(\m,\f)|\geq c|\m-\m_0(\f)|$ for $\m\in [0,1]$ and $\f\in [0,1]$. Then we have
\begin{align*}
|\pa_{\m}S_{\f}(\m)|=\m_0(\f)^{-1}\cdot |(\pa_{\m}S)(\m_0(\f)\m,\f)|\geq c|\m-1|
\end{align*}
for $\m\in \supp \c_{\f}(\cdot,\l)$ and $\f\in (0,1]$.
Moreover, for $\a\in\mathbb{N}$ and $\a'\in \mathbb{N}\setminus\{0\}$, we have
\begin{gather*}
|\pa_{\m}^{\a}\c_{\f}(\m,\l)|=\l^{\n} \m_0(\f)^{\a}|(\pa_{\m}^{\a}\c)(\m_0(\f)\m,\l,\f)|\\
\phantom{|\pa_{\m}^{\a}\c_{\f}(\m,\l)|}{}\lesssim \l^{\n}\m_0(\f)^{\a} (\m_0(\f) \m)^{2\n-\a}\bigl(1+\l \m_0(\f)^2 \m\bigr)^{-\n}\lesssim \m^{\n-\a},\\
\supp (\c_{\f}(\cdot,\l))\subset \big\{\m_0(\f)^{-1}\l^{-\frac{1}{2}}\leq \m\leq \m_0(\f)^{-1}\big\},\qquad \big|\pa_{\m}^{1+\a'}S_{\f}(\m)\big|\lesssim |\m_0(\f)|^{\a'-1}
\end{gather*}
Thus the stationary phase theorem (see Lemma \ref{stphaselem} with $k=2$, $j=0$ and $x_0=1$) implies
\begin{align*}
\m_0(\f)\l^{-\n}|I_2(\l,\f)|
\lesssim \m_0(\f)\l^{-\n}\cdot \bigl(\l \m_0(\f)^2\bigr)^{-\frac{1}{2}}=\l^{-\frac{1}{2}-\n},
\end{align*}
where we use $\l \m_0(\f)^2\gtrsim 1$.

Next, we consider the estimate of $I_{1}(\l,\f)$. We note that $|\pa_{\m}S_{\f}(\m)|\gtrsim 1$ for $\m\in \supp \g_1$.
Then the non-stationary phase theorem (see Lemma \ref{singularasym} with $\m=\n$) implies
\begin{align*}
\m_0(\f)\l^{-\n}|I_1(\l,\f)|\lesssim \m_0(\f)\l^{-\n}\cdot \bigl(\l \m_0(\f)^2\bigr)^{-\n-1}\lesssim \f\l^{-\n}\cdot \bigl(\l \f^2\bigr)^{-\frac{1}{2}}=\l^{-\frac{1}{2}-\n},
\end{align*}
where we use $\l \m_0(\f)^2\gtrsim 1$.

Finally, we consider the term $I_{3}(\l,\f)$. We note that $|\pa_{\m}S_{\f}(\m)|\gtrsim (|\m|+1)$ for $\m\in \supp \g_3$. Then, by integrating by parts $N(>1)$ times, we have
\begin{align*}
\m_0(\f)\l^{-\n}|I_3(\l,\f)|\lesssim \m_0(\f)\l^{-\n}\cdot \bigl(\l \m_0(\f)^2\bigr)^{-N}\lesssim \m_0(\f)\l^{-\n}\cdot \bigl(\l \m_0(\f)^2\bigr)^{-\frac{1}{2}}=\l^{-\frac{1}{2}-\n},
\end{align*}
where we use $\l \m_0(\f)^2\gtrsim 1$. This completes the proof.
\end{proof}

\appendix\setcounter{section}{0}
% Original source lines 1539--1612
\section{Asymptotic behavior of special functions}

\subsection{Leibniz's rule}

We frequently use the following formula, which directly follows from Leibniz's rule.

\begin{Lemma}\label{Leibnizrule}
Let $r$ be a smooth function and $N\in\mathbb{N}\setminus\{0\}$. Then there exist smooth functions~$b_{jN}$ such that
\begin{gather*}
\bigl(\pa_x\circ r(x)^{-1}\bigr)^N=\frac{1}{r(x)^N}\!\sum_{j=0}^Nb_{jN}(x)\pa_{x}^j\qquad \text{with}\quad |b_{jN}(x)|\lesssim \sum_{k=1}^{N-j}\sum_{\substack{\ell_1+\cdots+\ell_k=N-j,\\ \ell_1,\hdots,\ell_k\geq 1}}\prod_{i=1}^k\!\frac{\big|r^{(\ell_i)}(x)\big|}{|r(x)|}.
\end{gather*}
\end{Lemma}

\begin{Remark}
More precisely, we have
\begin{align*}
\bigl(\pa_x\circ r(x)^{-1}\bigr)^N=\frac{1}{r(x)^N}\pa_x^N+\frac{1}{r(x)^N}\sum_{j=0}^{N-1}\Bigg(\sum_{k=1}^{N-j}\sum_{\substack{\ell_1+\cdots+\ell_k=N-j,\\ \ell_1,\hdots,\ell_k\geq 1}}C_{\ell_1\ell_2\hdots \ell_k}^{jN}\prod_{i=1}^{k}\frac{r^{(\ell_i)}(x)}{r(x)} \Bigg)\pa_{x}^j
\end{align*}
with constants \smash{$C_{\ell_1\ell_2\hdots \ell_m}^{jN}>0$}.
\end{Remark}

\subsection{Non-stationary phase theorem with a singular amplitude}

The following lemma is more or less well known and an easy consequence of integration by parts. We give a proof for completeness of the paper.

\begin{Lemma}\label{singularasym}
Let $\m>-1$, $\c\in C^{\infty}((0,\infty))$ which is supported in $x\leq 10$ such that for $\a\in\mathbb{N}$, we have $|\pa_x^{\a}\c(x)|\lesssim |x|^{\m-\a}$ for $x\in (0,\infty)$. Let $f\in C^{\infty}(\R)$ satisfy $\im f(x)\geq 0$, $f'(x)\neq 0$ for~${x\in \supp \c}$ and $f(0)\in \R$. We define
\begin{align*}
b_{\m}(\l):={\rm e}^{-{\rm i}\l f(0)}\int_{0}^{\infty}\c(x){\rm e}^{{\rm i}\l f(x)}{\rm d}x\qquad \text{for}\quad \l\gtrsim 1.
\end{align*}
Then for each $\a\in\mathbb{N}$
\begin{align}\label{singbmues}
|\partial_{\l}^{\a}b_{\m}(\l)|\lesssim \l^{-\m-1-\a},\qquad \l\gtrsim 1.
\end{align}
In particular, $\big|\int_{0}^{\infty}\c(x){\rm e}^{{\rm i}\l f(x)}{\rm d}x\big|\lesssim \l^{-\m-1}$ for $\l\gtrsim 1$.
\end{Lemma}

\begin{proof}
We only need to prove these estimates for sufficiently large $\l$.

First, we prove \eqref{singbmues} for $\a=0$. Let $\chi\in C_c^{\infty}(\R)$ such that $\chi(x)=1$ for $|x|\leq 1/2$ and~${\chi(x)=0}$ for $|x|\geq 1$. Define $\chi_{\l}(x)=\chi(\l x)$ and $\overline{\chi_\l}(x)=1-\chi_{\l}(x)$. Then we have
\begin{align}\label{singbmues1}
\left|\int_{0}^{\infty}\c(x)\chi_{\l}(x){\rm e}^{{\rm i}\l f(x)}{\rm d}x\right|\lesssim \int_{0}^{1/\l}x^{\m}{\rm d}x\lesssim {\l}^{-\m-1}.
\end{align}
On the other hand, the integration by parts yields
\begin{align*}
\left|\int_{0}^{\infty}\c(x)\overline{\chi_{\l}}(x){\rm e}^{{\rm i}\l f(x)}{\rm d}x\right|=\l^{-N}\left|\int_{0}^{\infty}L^N(\c(x)\overline{\chi_{\l}}(x)){\rm e}^{{\rm i}\l f(x)}{\rm d}x\right|,
\end{align*}
where we define $L=\pa_x\circ ({\rm i}f'(x))^{-1}$. By using Leibniz's rule, Lemma \ref{Leibnizrule} and the assumption $f'\neq 0$ on $\supp \c$, we obtain $|L^{N}(\c(x)\overline{\chi_\l}(x))|\lesssim |x|^{\m-N}$. Hence, for $N>\m+1$,
\begin{align}\label{singbmues2}
\left|\int_{0}^{\infty}\c(x)\overline{\chi_{\l}}(x){\rm e}^{{\rm i}\l f(x)}{\rm d}x\right|\lesssim \l^{-N}\int_{\frac{1}{2\l}}^{10}x^{\m-N}{\rm d}x \lesssim \l^{-\m-1}.
\end{align}
The inequalities \eqref{singbmues1} and \eqref{singbmues2} imply \eqref{singbmues} for $\a=0$.

To deal with the case $\a\geq 1$, we write
\begin{align*}
\pa_{\l}^{\a}b_{\m}(\l)={\rm i}^{\a}{\rm e}^{-{\rm i}\l f(0)}\int_0^{\infty}(f(x)-f(0))^{\a}\c(x){\rm e}^{{\rm i}\l f(x)}{\rm d}x.
\end{align*}
By Taylor's theorem, we have $\pa_x^{N}(\c(x)(f(x)-f(0))^{\a})=O\bigl(|x|^{\m+\a-N}\bigr)$ for $x\in \supp \c$. Moreover, $\big|{\rm e}^{-{\rm i}\l f(0)}\big|=1$ due to the assumption $f(0)\in \R$.
Hence the same argument as the case $\a=0$ gives \eqref{singbmues} for $\a\geq 1$.
\end{proof}

\subsection{Stationary phase theorem}

In this paper, we use the following versions of the stationary phase theorem (the van der Corput lemma). The proof is same as the proof of \cite[Lemma 1.1.2]{So} (see also in \cite[p.~334]{S}). The statement about uniformity in a parameter follows from its proof and Lemma \ref{Leibnizrule}. We consider an integral
\begin{align*}
I(\l):=\int_{\R}\c(x){\rm e}^{{\rm i}\l f(x)}{\rm d}x\qquad \text{for}\quad \l\gtrsim 1.
\end{align*}

\begin{Lemma}\label{stphaselem}
Let $k\geq 2$ be an integer and $j\geq 0$, $x_0\in \R$, $\c\in C^{\infty}(\R\setminus \{x_0\})\cap C_c(\R)$ and~${f\in C^{\infty}(\R;\R)}$ such that $|\pa_x^{\a}\c(x)|\lesssim |x-x_0|^{j-\a}$, $|f(x)|\lesssim |x-x_0|^k$ and $|f'(x)|\gtrsim |x-x_0|^{k-1}$ for $x\in \supp \c\setminus \{x_0\}$.
Then \smash{$|I(\l)|\lesssim \l^{-\frac{j+1}{k}}$} for $\l\gtrsim 1$.
\end{Lemma}

\begin{thebibliography}{99}
\bibitem[27]{So}
Sogge C.D., Fourier integrals in classical analysis, \textit{Cambridge Tracts
 in Math.}, Vol. 105, \href{https://doi.org/10.1017/CBO9780511530029}{Cambridge University Press}, Cambridge, 1993.

\bibitem[28]{S}
Stein E.M., Harmonic analysis: real-variable methods, orthogonality, and
 oscillatory integrals, \textit{Princeton Math. Ser.}, Vol.~43, Princeton
 University Press, Princeton, NJ, 1993.

\end{thebibliography}
\end{document}
